\documentclass{article}
\usepackage[utf8]{inputenc}
\usepackage{a4wide}

\begin{document}

\section*{Box  {CT}}

{ในแผนที่ขนาด\texttt{N\ แถว\ x\ M}คอลัมน์ที่มีทั้งช่องว่างและช่องที่มีสิ่งกีดขวาง

คุณต้องการเคลื่อนย้ายกล่องขนาด\texttt{2x2}ให้ผ่านแผนที่ดังกล่าวจากด้านบนสุดให้ถึงด้านล่าง

คุณสามารถเริ่มต้นจากใดก็ได้ในขอบด้านบน โดยที่กล่องต้องอยู่ในขอบเขตของแผนที่ทั้งหมด

ในการเคลื่อนย้ายกล่องนั้นคุณไม่สามารถเคลื่อนออกจากพื้นที่นี้ได้

และ\textbf{กล่องสามารถเคลื่อนที่ไปได้ในทิศทางขึ้น ลง ซ้าย หรือขวาเท่านั้น

ไม่สามารถเคลื่อนที่แนวทแยงได้}}\\



{{ตัวอย่างการเคลื่อนที่แสดงในภาพด้านล่าง ด้านซ้ายเป็นตัวอย่างแผนที่

จุด\texttt{(.)}แทนช่องว่าง และ ชาร์ป\texttt{(\#)}แทนช่องที่มีสิ่งกีดขวาง}\\

}



{{\includegraphics[width=6.33333in,height=2.08333in,alt={image.png}]{/public/upload/bd96c50f9e.png}\\

}}



{{คุณสามารถเคลื่อนกล่องจากด้านบนไปด้านล่าง ผ่านทางช่องที่ทำเครื่องหมาย x ไว้ในด้านขวา\\

}}



\subsection*{Input}
{บรรทัดแรกมีจำนวนเต็ม\texttt{N}และ\texttt{M}\texttt{(2\ \textless{}=\ N\ \textless{}=\ 30;\ 2\ \textless{}=\ M\ \textless{}=\ 30)}จากนั้นอีก\texttt{N}บรรทัดจะมีสตริงความยาว\texttt{M}ตัวอักษร

แต่ละตัวอักษรแทนช่องในแผนที่

โดยอาจจะเป็นสิ่งกีดขวาง\texttt{(\textquotesingle{}\#\textquotesingle{})}หรือช่องว่าง\texttt{(\textquotesingle{}.\textquotesingle{})}}\\



\subsection*{Output}
{โปรแกรมคุณควรเขียนผลลัพธ์ในบรรทัดเดียว

โดยแสดง\texttt{yes}หากสามารถพากล่องผ่านแผนที่นี้ได้

และ\texttt{no}หากไม่สามารถทำได้}\\



\end{document}
